% Шаблон отчёта make_report. Компиляция: tectonic report.tex
\documentclass[11pt,a4paper]{article}

\usepackage{fontspec}
\usepackage{polyglossia}
\setmainlanguage{russian}
\setotherlanguage{english}
\setmainfont{Times New Roman}
\newfontfamily\cyrillicfont{Times New Roman}
\setsansfont{Helvetica}
\newfontfamily\cyrillicfontsf{Helvetica}
\setmonofont{Menlo}[Scale=0.9]
\newfontfamily\cyrillicfonttt{Menlo}[Scale=0.9]

\usepackage[margin=2cm]{geometry}
\usepackage{amsmath,amssymb}
\usepackage{booktabs,array,multirow,tabularx,longtable}
\usepackage{graphicx}
\usepackage{float}
\usepackage[font=small,labelsep=period]{caption}
\usepackage{subcaption}
\usepackage[shortlabels]{enumitem}
\usepackage{xcolor}
\usepackage[colorlinks=true,linkcolor=black,urlcolor=blue!50!black,citecolor=black]{hyperref}

\setlist{nosep,leftmargin=1.2em}
\setlength{\parindent}{0pt}
\setlength{\parskip}{4pt}
\graphicspath{{figures/}}

% Роли цветов (совпадают с фигурами)
\definecolor{ml}{HTML}{2A78D6}      % наша модель
\definecolor{ppub}{HTML}{EB6834}    % Phosphoros (опубликованный)
\definecolor{pfree}{HTML}{1BAF7A}   % Phosphoros без информации N
\definecolor{poff}{HTML}{EDA100}    % официальный PHZ

\newcommand{\NMAD}{\mathrm{NMAD}}
\newcommand{\dz}{\Delta z/(1+z)}
\newcommand{\fnref}[1]{\textsuperscript{\ref{#1}}}  % повторная ссылка на сноску \footnote{\label{fn:key}...}
\newcolumntype{Y}{>{\raggedright\arraybackslash}X}

\title{Отчёт за <<ДАТА>>}
\author{}
\date{}

\begin{document}
\maketitle

\section{Сетап эксперимента}

\subsection{Данные}

\begin{table}[H]\centering\small\caption{Входы модели.}
\begin{tabularx}{\linewidth}{llY}
\toprule Вход & Форма & Содержимое \\ \midrule
Изображение & <<C$\times$H$\times$W>> & <<слои>> \\
Таблица & <<N$\times$K>> & <<полосы, столбцы>> \\
\bottomrule\end{tabularx}\end{table}

\begin{table}[H]\centering\small\caption{Объём выборок.}
\begin{tabular}{lrrrr}
\toprule Этап & F & S & N & Всего \\ \midrule
<<этап>> & <<…>> & <<…>> & <<…>> & <<…>> \\
\bottomrule\end{tabular}\end{table}

\begin{table}[H]
\centering
\caption{Обзор меток}
\begin{tabular}{llrl}
\toprule
Обзор & Поле & Число & Роль \\
\midrule
<<обзор>> & <<поле>> & <<число>> & <<роль>> \\
\bottomrule
\end{tabular}
\end{table}

% \begin{figure}[H]\centering\includegraphics[width=\linewidth]{fig_label_dist.pdf}\caption{}\end{figure}  % подписи к рисункам пустые

\subsection{Трейн-выборка}
\begin{itemize}
  \item <<поля, метки, немеченые данные, синтетика, перевод/нормировка>>
\end{itemize}

\subsection{Тест-выборка}
\begin{itemize}
  \item <<наборы поля N, отбор, исключения, подтверждающий / описательный>>
\end{itemize}

\subsection{Модели}
\paragraph{CLIP (предобучение).}
\begin{itemize}
  \item <<энкодеры, размерности, лосс, температура, данные, эпохи, оптимизатор, батч, аугментации, параметры>>
\end{itemize}
\paragraph{DL-головы и дообучение.}
\begin{itemize}
  \item <<архитектура, лосс, гиперпараметры, сиды, ансамбль>>
\end{itemize}
\paragraph{Базовые методы.}
\begin{itemize}
  \item <<Phosphoros, официальный PHZ: что это и настройки>>
\end{itemize}

\section{Результаты эксперимента}

\subsection{Метрики}

\begin{table}[H]
\centering
\caption{<<таблица метрик>>}
\begin{tabular}{lrrr}
\toprule
Набор & <<метод 1>> & <<метод 2>> & $\Delta$ [95\% ДИ] \\
\midrule
<<набор>> & <<NMAD / выбросы>> & <<…>> & <<…>> \\
\bottomrule
\end{tabular}
\end{table}

\subsection{Статистические тесты}
\begin{itemize}
  \item <<бутстрэп: что ресэмплируется, B, seed; решение по предрегистрации>>
\end{itemize}

\subsection{Графики}
% рассеяние, моды по бинам + ломаная, NMAD по бинам z и VIS, forest-plot ΔNMAD

\end{document}
